\begin{lemma}[Transporting a hierarchy base across a same-color presentation]
Fix an extended Tashkinov tree \(\mathcal E\), generated by a coloring
\(\varphi\), and a \noderef{HierarchyBase} \(B\) for
\((\mathcal E,\varphi)\).  Let
\[
 C:\noderef{CorrespondingExtendedTashkinovTree}
       (\mathcal E,\varphi,U)
\]
be any corresponding presentation with the same generating coloring and
with an arbitrary exact displayed tree \(U\).  Then the replacement ETT
\(\mathcal E'=C.\mathit{replacementETT}\) has a hierarchy base \(B'\)
under \(\varphi\) such that
\[
 B'.\mathit{core}=B.\mathit{core},\qquad
 B'.\mathit{closure}=B.\mathit{closure},\qquad
 B'.\mathit{supportPosition}=B.\mathit{supportPosition}
\]
in the last equality at the level of numerical index values.
\end{lemma}
\begin{proof}
Write \(n=\mathcal E.\mathit{rounds}\), let \(T_n\) be the terminal ladder
tree of \(\mathcal E\), and let \(T'_n\) be the terminal ladder tree in the
replacement series displayed by \(C\).  The positivity field of
\noderef{ExtendedTashkinovTree} gives \(n\geq1\), so \(q=n-1\) is a valid
positive-stage index and \(q+1\) is the terminal rung.  The rung agreement,
series presentation, and ladder-tree agreement fields of
\noderef{CorrespondingExtendedTashkinovTree}, evaluated at \(q\), give an
equality of exact ordered edge-valued tree-sequences
\[
                         T'_n=T_n.                              \tag{1}
\]
The same correspondence definition, now at the terminal stored index,
gives
\[
 \Theta'_n=\Theta_n,qquad F'_n=F_n,                            \tag{2}
\]
where the second equality is equality of sets of exact connecting-edge
identities.

We first establish the only item not explicitly stored in the
correspondence certificate:
\[
 \Theta_n=PE\quad\Longrightarrow\quad v'_n=v_n,                 \tag{3}
\]
where \(v_n,v'_n\) are the terminal chosen vertices.  Assume
\(\Theta_n=PE\).  By (2), the replacement terminal kind is also \(PE\).
The certified-step field of \noderef{TashkinovSeries}, together with the
constructors in \noderef{TashkinovSeriesStep}, shows that both terminal
certificates must therefore be their \(PE\) constructors: the initial,
SE, and RE constructors have respectively kind none, SE, and RE.

Unpack the two PE constructors, and call their
\noderef{MaximumDefectiveChoice}s \(M\) and \(M'\).  Their connecting-edge
fields say
\[
                         F_n=\{M.\mathit{edge}\},
 \qquad F'_n=\{M'.\mathit{edge}\}.
\]
Together with (2), singleton equality identifies the exact edge identities
\[
                         M'.\mathit{edge}=M.\mathit{edge}.       \tag{4}
\]
This is equality of multigraph edges, not merely equality of endpoint
pairs.

By the definition of \noderef{MaximumDefectiveChoice}, the endpoint
\(M.\mathit{vertex}\) lies in \(V(T_n)\), its other endpoint lies outside
\(V(T_n)\), and the ends of its selected edge are exactly the unordered
pair of those two vertices.  The analogous assertions hold for \(M'\) and
\(T'_n\).  Use (1) to regard both inside endpoints as lying in the same
set \(V(T_n)\), and both other endpoints as lying outside it.  Equations
(1) and (4) identify the two unordered endpoint pairs.  Equality of
unordered pairs has two orientations.  The orientation matching first
endpoint to first endpoint yields
\(M'.\mathit{vertex}=M.\mathit{vertex}\).  The crossed orientation would
identify the inside endpoint of one choice with the outside endpoint of
the other, contradicting their membership and nonmembership in the common
set \(V(T_n)\).  Hence only the first orientation is possible.  Finally,
the chosen-vertex fields of the two PE constructors identify \(v_n\) and
\(v'_n\) with these two equal inside endpoints.  This proves (3).

We now construct \(B'\).  Retain the exact trees
\[
 B'.\mathit{core}:=B.\mathit{core},\qquad
 B'.\mathit{closure}:=B.\mathit{closure}.                       \tag{5}
\]
Transport the support position of \(B\) across the length equality induced
by (1).  This produces an index in \(T'_n\) with the same numerical value;
its positivity follows from the support-positivity field of
\noderef{HierarchyBase}.  Since both the old and new hierarchy bases use
the same coloring \(\varphi\) as their ETT generating parameter, their
terminal missing-color requirement is the reflexive equality
\(\overline\varphi(T_n)=\overline\varphi(T_n)\), after (1).

Suppose first that \(\Theta'_n=PE\).  Equation (2) gives
\(\Theta_n=PE\), so the parallel clause of the old
\noderef{HierarchyBase} supplies three facts: the recorded old chosen
vertex is the vertex of \(T_n\) at its support position, the retained tree
\(B.\mathit{closure}\) is a
\noderef{TashkinovClosure} of the corresponding
\noderef{TreeSequencePrefix}, and \(B.\mathit{core}\) carries an
\noderef{OrderedTreeUnion} of \(T_n\) with that closure.  Equality (3),
the exact tree equality (1), and the transported support position prove
the required chosen-vertex equation for \(B'\).  Equality (1), equality
of the numerical support positions, and proof irrelevance identify the
two exact prefix structures.  Thus the old closure certificate is exactly
the required closure certificate for \(B'\).  Substituting (1) and (5) in
the old ordered-union certificate gives the required ordered union for
\(B'\).

Suppose instead that \(\Theta'_n\ne PE\).  Equation (2) implies
\(\Theta_n\ne PE\).  The nonparallel clause of the old
\noderef{HierarchyBase} gives \(B.\mathit{core}=T_n\).  Combining this
with (1) and (5) yields \(B'.\mathit{core}=T'_n\), which is precisely the
new nonparallel clause.

All fields of a hierarchy base have now been verified.  By (5) its core
and closure are literally the old ones, and by construction its support
position has the same numerical value, proving all three asserted output
equalities.
\end{proof}
